\loai{Bài toán bọt khí nổi lên từ đáy hồ}
\begin{vidu} %[1P5-2.3G][Loai3]
	Một bọt khí nhỏ, nổi lên từ đáy hồ, khi đến mặt nước thì thể tích của nó tăng gấp $1,3$ lần. Tính độ sâu của đáy hồ, biết trọng lượng riêng của nước là $d=10^4\ N/m^3$, áp suất khí quyển là $p_0=10^5\ N/m^2$. Coi nhiệt độ của nước ở các điểm là như nhau.
	\choice
	{1 m}
	{2 m}
	{\True 3 m}
	{4 m}
	\loigiai{
		immini[thm]{		
			\begin{itemize}
				\item Ở đáy hồ: $p=p_0+dh;\ V$
				\item Ở mặt hồ: $p_0;\ V_0=1,3V$
				\item Do nhiệt độ không đổi, theo định luật Bôi-lơ - Ma-ri-ốt:
				\begin{align*}
					p_0V_0=pV\Rightarrow (p_0+dh).V=1,3Vp_0\Rightarrow  h=\dfrac{0,3p_0}{d} =\dfrac{0,3.10^5}{10^4}=3\ m.
				\end{align*}
				\begin{note}
					Khi bọt khí đến mặt nước thì áp suất của nó chính là áp suất khí quyển nên thể tích của nó không thay đổi nữa.\\
					Công thức Pascal: $p_1=p_0+Dgh=p_0+dh$.
				\end{note}
			\end{itemize}
		}
		{\begin{tikzpicture}[scale=1,thick,>=stealth']
				\fill[green!20!cyan](0,0)--(6,0)--(6,4)--(0,4)--cycle;
				\fill[pattern=north east lines] (0,0)rectangle(6,-0.1);
				\draw(0,0)--(6,0)(0,4)--(6,4);
				\tkzLabelPoint[below left](o){$O$}
				\draw[<->](0.3,0)--node[midway,right]{$h$}(0.3,4);
				\draw[fill=white](1.5,0.3) circle (0.3cm)node{$V_1$};
				\draw[fill=white](1.5,3.5) circle (0.5cm)node{$V_2$};
				\draw[->](2.5,0.5)node[above right]{$p_1=p_0+Dgh$}--(2.5,0);
				\draw[->](2.5,4.5)node[above right]{$p_0=10^5pa$}--(2.5,4);
		\end{tikzpicture}}
	}
\end{vidu}

\begin{vidu} %[1P5-2.3K][Loai3]
	Một bong bóng khí ở độ sâu 5 m có thể tích thay đổi như thế nào khi nổi lên mặt nước. Cho áp suất tại mặt nước là $10^5\ Pa$, khối lượng riêng của nước là $1000\ kg/m^3$, gia tốc trọng trường là $10\ m/s^2$.
	\choice
	{Giảm 3 lần}
	{Giảm 1,5 lần}
	{Tăng 3 lần}
	{\True Tăng 1,5 lần}
	\loigiai{
		\begin{itemize}
			\item Trạng thái 1, ở độ sâu 5 m: 
			\begin{align*}
				V_1,\ p_1=\rho gh + 10^5= 1000.10.5 +10^5= 1, 5.10^5\ Pa.
			\end{align*}
			\item Trạng thái 2, ở mặt nước: $V_2$, $p_2=10^5\ Pa$.
			\item Theo định luật Bôi-lơ - Ma-ri-ốt: 
			\begin{align*}
				p_1.V_1=p_2.V_2\Rightarrow \dfrac{V_2}{V_1} = \dfrac{p_1}{p_2} = \dfrac{1, 5.10^5}{10^5} = 1,5 \Rightarrow\ \ct{Thể tích tăng 1,5 lần}
			\end{align*}
		\end{itemize}
	}
\end{vidu}
\loai{Bài toán chất lỏng trong ống thủy tinh}
\begin{vidu}%[1P5-2.3T][Loai4]
	immini[thm]{
		Một ống thủy tinh hình trụ tiết diện S nhỏ, một đầu kín một đầu hở. Bên trong ống thủy tinh chứa khí lí tưởng được ngăn với bên ngoài bởi cột thủy ngân có chiều dài 15 cm. Tại thời điểm ban đầu ống để thẳng đứng đầu hở quay lên trên (hình vẽ), người ta đo được chiều dài của cột không khí bên trong ống là 30 cm. Biết áp suất khí quyển là $1,013.10^5$ Pa, lấy $g=10\ m/s^2$, khối lượng riêng của thủy ngân là 13600 $kg/m^3$. Coi quá trình là đẳng nhiệt. Tính chiều dài của cột không khí trong các trường hợp sau (lưu ý giọt thủy ngân không bị rơi khỏi ống thủy tinh).
		\begin{listEX}[2]
			\item Ống được đặt nằm ngang 
			\item Ống được đặt thẳng đứng, miệng ống ở dưới 
	\end{listEX}}
	{
		\begin{tikzpicture}[scale=1,thick,>=stealth']
			\draw[fill=gray!40,opacity=0.5](0,0)node[below left,opacity=1]{$\ell_0$}--++(270:0.5)--++(0:1)--++(90:0.5)--cycle;
			\draw(0,1)--++(270:4)--++(0:1)--++(90:4);
			\draw[<->] (1.3,-0.5)--node[midway,right]{$\ell_1$}(1.3,-3);
		\end{tikzpicture}
	}
	
	%hinh1
	\loigiai{
		\textbf{Phân tích bài toán}:
		\begin{itemize}
			\item Để giải bài toán cần sử dụng công thức tính thể tích hình trụ $V=Sh$ với S là tiết diện ngang, h là chiều cao của ống.
			\item Công thức tính áp suất $p=\dfrac{F}{S}$ với $F=mg=\rho Vg=\rho \ell_0Sg$ là trọng lượng của vật nén lên diện tích S.\\
			$\ell_0=15\ cm=0,15\ m;\  \ell_1=30\ cm = 0,3\ m;\ p_0=1,013.10^5\ Pa$
			\item Trạng thái 1: Khi ống thẳng đứng, miệng ống ở trên: $V_1=\ell_1S;\ p_1=p_0+\rho \ell_0g$
		\end{itemize}
		a) Ống đặt nằm ngang ở trạng thái 2: $V_2=\ell_2S;\ p_2=p_0$
		\begin{align*}
			V_2p_2=V_1p_1 \Rightarrow \ell_2=0,36\ m
		\end{align*}
		b) Ống đặt thẳng đứng miệng ống quay xuống dưới ở trạng thái 3: $V_3=\ell_3S;\ p_3=p_0-\rho \ell_0g$
		\begin{align*}
			V_3p_3=V_1p_1 \Rightarrow \ell_3=0,45\ m
		\end{align*}
	\begin{note}
		\begin{itemize}
			\item \textbf{Nguyên lý Pascal} hay định luật Pascal là độ tăng áp suất lên một chất lỏng chứa trong bình kín được truyền nguyên vẹn cho mọi điểm của chất lỏng và thành bình do nhà bác học người Pháp Blaise Pascal phát hiện khi lợi dụng tính chất khó nén của nước và nguyên lý bình thông nhau.
			\begin{align*}
				p=p_a + \rho gh 
			\end{align*}
			Trong đó: 
			\begin{itemize}
				\item $p_a$ là áp suất từ bên ngoài nén lên mặt chất lỏng.
				\item $\rho$ là khối lượng riêng của chất lỏng. 
				\item h là khoảng cách tính từ bề mặt chất lỏng đến điểm cần xét.
			\end{itemize}
			\item Nếu áp suất trong bài tính theo mmHg thì áp suất trong ống được tính theo công thức: $p=p_0 \pm \ell_0$.
		\end{itemize}
	\end{note}
}
\end{vidu}
\loai{Bài toán lực tác dụng khi có sự chênh lệch áp suất}
\begin{vidu} %[1P5-2.3G][Loai 5]
Một lượng không khí có thể tích $240\ cm^3$ bị giam trong một xi lanh đặt nằm ngang có pit-tông đóng kín, diện tích của pit-tông là $24\ cm^2$. Áp suất khí trong xi lanh bằng áp suất ngoài là 100 kPa. Cần một lực bằng bao nhiêu để di chuyển pit-tông một đoạn 2 cm theo chiều nén khí? Bỏ qua ma sát giữa pit-tông và thành xi lanh. Coi quá trình xảy ra là đẳng nhiệt.
\choice
{\True 60 N}
{40 N}
{20 N}
{80 N}
\loigiai{
	\Binhluan{
		\begin{itemize}
			\item Lúc đầu áp suất và thể tích khí trong xi lanh là $p_0$, $V_0$.
			\item Sau khi đẩy pit-tông 2 cm theo chiều nén khí thì áp suất là $p_0+ \dfrac{F}{s},\ V_0- S.\ell$.
			\item  Theo định luật Bôi-lơ - Ma-ri-ốt: 
			\begin{align*}
				&p_0V_0=\left(p_0 + \dfrac{F}{S}\right)\left(V_0 - S\ell\right)\\
				&p_0S\ell = F\left(\dfrac{V_0}{S} - \ell\right)\\
				\Rightarrow\ &F =\dfrac{p_0S^2\ell}{V_0 - \ell S}=\dfrac{100\cdot 10^3\cdot\left(24.10^{ - 4}\right)^2\cdot 2\cdot 10^{ - 2}}{240.10^{ - 6} - 2.10^{ - 2}\cdot 24\cdot 10^{ - 4}}= 60\ N.
			\end{align*}	
		\end{itemize}
	}
	{Khi ấn xi lanh ta gây thêm cho khối khí bên trong áp suất $p=\dfrac{F}{S}$. Tính thêm áp suất khí quyển $p_0$ nên khi pit-tông cân bằng thì 
		\begin{align*}
			p_2=p_0+\dfrac{F}{S}
		\end{align*}	
	}
}
\end{vidu}
